\documentclass[a4paper,10pt]{article}

\usepackage[utf8]{inputenc}
\usepackage[T1]{fontenc}
\usepackage{geometry}
\usepackage{booktabs}
\usepackage{array}
\usepackage{tabularx}
\usepackage{longtable}
\usepackage{listings}
\usepackage{xcolor}
\usepackage[hidelinks]{hyperref}
\usepackage{amsmath}

\geometry{
    a4paper,
    top=1.8cm,
    bottom=1.8cm,
    left=1.8cm,
    right=1.8cm
}

\setlength{\parindent}{0pt}
\setlength{\parskip}{5pt}
\renewcommand{\arraystretch}{1.18}

\lstdefinestyle{cstyle}{
    language={[ANSI]C},
    basicstyle=\ttfamily\small,
    keywordstyle=\bfseries,
    commentstyle=\itshape,
    numbers=left,
    numberstyle=\tiny,
    numbersep=7pt,
    frame=single,
    breaklines=true,
    showstringspaces=false,
    tabsize=4
}

\newcommand{\hex}[1]{\texttt{0x#1}}
\newcommand{\code}[1]{\texttt{#1}}

\title{\textbf{DEOS CAN-FD Transport Adapter ICD}\\
\large Sürüm 1.1 -- Fiziksel CAN-FD Taşıma Katmanı}
\author{\textbf{Emir Arkalı} (\href{mailto:emirarkali07@gmail.com}{emirarkali07@gmail.com})\\DEOS -- Dokuz Eylül Otonom Sistemler}
\date{İlk Yayın: 28 Eylül 2026 \quad | \quad Son Revizyon: 13 Ekim 2026}

\begin{document}
\maketitle

\begin{center}
\small
\begin{tabularx}{\textwidth}{ll lX}
\toprule
\textbf{Doküman Tipi:} & Interface Control Document (ICD) & \textbf{Doküman Kodu:} & \code{DOC-ICD-TRP-CANFD-005} \\
\textbf{Doküman Sürümü:} & Sürüm 1.1 & \textbf{Protokol Sürümü:} & v1.1 (\hex{11}) -- Medya Bağımsız \\
\textbf{Yazar:} & Emir Arkalı (\texttt{emirarkali07@gmail.com}) & \textbf{Son Değiştiren:} & Emir Arkalı \\
\textbf{İlk Yayın Tarihi:} & 28 Eylül 2026 & \textbf{Son Revizyon Tarihi:} & 13 Ekim 2026 \\
\textbf{Durum:} & Yürürlükte / Onaylandı & \textbf{Taşıma / Kapsam:} & CAN-FD Taşıma Adaptörü (29-bit ID, DLC, BRS) \\
\bottomrule
\end{tabularx}
\end{center}

\subsection*{Doküman Revizyon Geçmişi}
\begin{center}
\small
\begin{tabularx}{\textwidth}{c c l X}
\toprule
\textbf{Revizyon} & \textbf{Tarih} & \textbf{Yazar / Değiştiren} & \textbf{Açıklama} \\
\midrule
v0.3 & 28 Eylül 2026 & Emir Arkalı & İlk taslak ve monolitik ICD v0.3 tanımı. \\
v1.0 & 04 Ekim 2026 & Emir Arkalı & Modüler CAN-FD taşıma adaptörü ayrımı. \\
v1.1 & 13 Ekim 2026 & Emir Arkalı & Sürüm 1.1 standardizasyonu, Evrensel Çerçeve dönüşüm mekanizması entegrasyonu ve kurumsal üstveri eklemesi. \\
\bottomrule
\end{tabularx}
\end{center}
\vspace{0.8em}

\section{Giriş ve Kapsam}

DEOS protokolü medya bağımsız (medya agnostik) bir mimariye sahiptir. Araç içi ana omurga haberleşmesi olarak ise \textbf{CAN-FD (Controller Area Network Flexible Data-Rate)} fiziksel katmanı ve veri bağı (Data Link) katmanı kullanılır.

Bu doküman, medyanın üstündeki soyut DEOS Evrensel Çerçevesi'nin (\code{Universal DEOS Transport Frame}) CAN-FD donanımı üzerinde nasıl temsil edildiğini, 29-bit Extended Identifier bit yerleşimini, CAN-FD DLC yuvarlama ve doldurma (padding) kurallarını, Bit Rate Switching (BRS) ve filtreleme kurallarını tanımlar.

\subsection{Evrensel Çerçeveden CAN-FD'ye Dönüşüm}

Çekirdek protokolde tanımlanan standart 9+N baytlık Evrensel Çerçeve (Universal Frame), CAN-FD omurgasına aktarılırken iki parçaya ayrıştırılır:
\begin{enumerate}
    \item \textbf{29-bit CAN Identifier (Arbitration Alanı):} Byte 0 (\code{Destination Node}), Byte 1 (\code{Source Node}), Byte 2 (\code{Priority}), Byte 4 (\code{Message Class}) ve Byte 5 (\code{Service ID}) alanları donanımsal bus arbitration ve filtreleme amacıyla 29-bit Extended ID'ye dönüştürülür.
    \item \textbf{CAN-FD Veri Alanı (Payload):} Byte 3 (\code{Version}), Byte 7 (\code{Sequence}), Byte 6 (\code{Command}), Byte 8 (\code{Payload Length}) ve ardından gelen dinamik $N$ bayt \code{Payload} doğrudan CAN-FD frame'inin veri alanına yazılır.
\end{enumerate}

\section{29-bit Extended CAN Identifier Yapısı}

DEOS protokolü standart (11-bit) yerine 29-bit Extended CAN Identifier kullanır. Standart DEOS mesajındaki yönlendirme ve öncelik alanları doğrudan 29-bitlik bu kimlik alanına haritalanır:

\begin{center}
\begin{tabular}{c c c c c}
\toprule
\textbf{Bitler} & \textbf{Alan Adı} & \textbf{Boyut} & \textbf{Açıklama} \\
\midrule
28--26 & \code{Priority} & 3 bit & Arbitration önceliği (\hex{0} en yüksek, \hex{7} en düşük) \\
25--22 & \code{Message Class} & 4 bit & Mesaj davranış sınıfı (\hex{1} COMMAND, \hex{2} RESPONSE ..) \\
21--16 & \code{Service ID} & 6 bit & Fonksiyonel servis kimliği (\hex{00} SYSTEM .. \hex{08} LIGHTING) \\
15--8  & \code{Destination Node} & 8 bit & Hedef düğüm adresi (\hex{00} -- \hex{FF}) \\
7--0   & \code{Source Node} & 8 bit & Kaynak düğüm adresi (\hex{00} -- \hex{FF}) \\
\bottomrule
\end{tabular}
\end{center}

Toplam uzunluk:
\[
3 + 4 + 6 + 8 + 8 = 29 \text{ bit}
\]

\subsection{CAN Identifier Kodlama Formülü}

C/C++ sürücü katmanında (\code{deos\_can\_id\_encode}) bu 29-bit kimlik şu şekilde bit kaydırmalarıyla hesaplanır:

\[
\text{CAN\_ID} = (P \ll 26) \;|\; (C \ll 22) \;|\; (S \ll 16) \;|\; (D \ll 8) \;|\; SRC
\]

Burada:
\begin{itemize}
    \item \(P\): \code{Priority} (3 bit),
    \item \(C\): \code{Message Class} (4 bit),
    \item \(S\): \code{Service ID} (6 bit),
    \item \(D\): \code{Destination Node} (8 bit),
    \item \(SRC\): \code{Source Node} (8 bit) değeridir.
\end{itemize}

\subsection{Arbitration ve Öncelik Davranışı}

CAN bus seviyesinde dominant bit (\code{0}), resesif bitten (\code{1}) üstündür. Bu sebeple sayısal olarak daha küçük olan CAN-ID hat üzerinde çatışmayı (arbitration) kazanır:
\begin{itemize}
    \item \code{Priority = 0} (\code{DEOS\_PRIO\_EMERGENCY}) en yüksek hat önceliğine sahiptir.
    \item \code{Priority = 5} (\code{DEOS\_PRIO\_NETWORK}) düşük hat önceliğine sahiptir.
\end{itemize}

\section{CAN-FD Veri Alanı Formatı (Data Field Layout)}

CAN-FD çerçevesinin veri alanında (Payload alanı) sabit 4 baytlık ortak başlık (\code{DEOS Common Header}) ve ardından komut verisi taşınır:

\begin{center}
\begin{tabular}{c c c l}
\toprule
\textbf{Byte Ofset} & \textbf{Alan Adı} & \textbf{Boyut} & \textbf{Açıklama} \\
\midrule
0 & \code{Version} & 1 Byte & Protokol sürümü (\hex{11} = v1.1) \\
1 & \code{Sequence} & 1 Byte & Düğüm bazlı artan mesaj takip sayacı \\
2 & \code{Command} & 1 Byte & İlgili servise ait komut kodu \\
3 & \code{Length} & 1 Byte & Arkadan gelen geçerli payload boyutu ($N$, 0--60 bayt) \\
\midrule
4 .. (4+$N$-1) & \code{Payload} & $N$ Byte & Komuta özel gerçek veri içeriği \\
\bottomrule
\end{tabular}
\end{center}

\section{CAN-FD DLC, Padding ve BRS Kuralları}

\subsection{DLC Yuvarlama (DLC Rounding)}

CAN-FD standardı 8 bayttan sonra doğrudan her bayt uzunluğunu desteklemez; veri boyutu yalnızca \textbf{12, 16, 20, 24, 32, 48 veya 64} bayt olabilir.

DEOS protokolünde:
\begin{itemize}
    \item \textbf{Length (Byte 3):} Yalnızca semantik payload bayt sayısını ifade eder (Örn: Far için 1 bayt).
    \item \textbf{DLC (Data Length Code):} Toplam semantik boyut ($4 + N$) en yakın geçerli CAN-FD boyutuna yuvarlanır.
    \item \textbf{Padding (Doldurma):} Yuvarlama nedeniyle arta kalan boş baytlar \hex{00} ile doldurulur ve alıcı yazılım tarafından dikkate alınmaz.
\end{itemize}

\subsection{Bit Rate Switching (BRS)}

Tüm DEOS CAN-FD paketlerinde:
\begin{itemize}
    \item \code{CAN\_FRAME\_IDE} (29-bit Extended ID),
    \item \code{CAN\_FRAME\_FDF} (CAN-FD Formatı),
    \item \code{CAN\_FRAME\_BRS} (Bit Rate Switch: Veri fazında yüksek hıza geçiş)
\end{itemize}
bayrakları aktiftir.

\section{Donanımsal Filtreleme (Hardware Filtering)}

Her düğümün CAN kontrolcüsü (STM32 FDCAN vb.), işlemciyi gereksiz kesmelerden (interrupt storm) korumak için donanımsal kabul filtreleri (acceptance filters) kullanır:
\begin{itemize}
    \item \textbf{Unicast Filtresi:} Destination = Kendi Node ID'si olan paketler kabul edilir.
    \item \textbf{CAN Broadcast Filtresi:} Destination = \hex{FE} (\code{DEOS\_NODE\_CAN\_BROADCAST}) olan paketler kabul edilir.
    \item \textbf{Global Broadcast Filtresi:} Destination = \hex{FF} (\code{DEOS\_NODE\_GLOBAL\_BROADCAST}) olan paketler kabul edilir.
    \item Diğer tüm paketler donanım seviyesinde sessizce elenir.
\end{itemize}

\end{document}
